\note{Moving to UC2, the authors used AdvNet for regression testing between BBR v1 and BBR v3, focusing on their fairness toward a competing CUBIC flow. In each candidate environment, they run two tests: first, CUBIC competing alongside BBR v1 (whose CUBIC throughput serves as the reference score), and second, CUBIC competing alongside BBR v3 (whose CUBIC throughput serves as the target score). To measure steady-state impact rather than transient jitter, AdvNet first varies bandwidth and latency, then holds the network conditions fixed for one minute, and measures throughput over the final ten seconds.\par \par In one adversarial environment discovered by AdvNet, when CUBIC competes with BBR v1, CUBIC achieves 12 times the throughput of BBR v1; yet when BBR v1 is replaced by BBR v3 under the exact same trace, the balance completely reverses, with BBR v3 achieving 11.6 times the throughput of CUBIC. In single-flow tests, AdvNet also found scenarios where BBR v1 achieves more than five times the throughput of BBR v3.\par \par Why do we see such dramatic reversals between versions? As the authors note, these outcomes are a direct consequence of the suboptimal BBR behaviors—such as rate collapse and slow convergence—uncovered in the earlier case studies. Even though BBR v3 was designed to improve fairness in typical settings, adversarial testing reveals specific dynamic conditions where version changes trigger severe regressions.
\par\smallskip Source: AdvNet paper §4.4.5; introduction}
